% Figure : diagramme de Feynman pour la diffusion Raman au premier ordre
% Compatible pdflatex.
% Préambule requis :
%   \usepackage{tikz}
%   \usetikzlibrary{decorations.pathmorphing}
%   \usetikzlibrary{decorations.markings}
%   \usetikzlibrary{arrows.meta}
%
% Style réutilisable pour une ligne de fermion avec flèche au milieu :
%   \tikzset{
%     fermion/.style={
%       postaction={decorate},
%       decoration={markings, mark=at position 0.5 with {\arrow{Stealth}}}
%     }
%   }

\begin{figure}[H]
\centering
\begin{tikzpicture}[
    line/.style={thick},
    >=Stealth,
    fermion/.style={
      postaction={decorate},
      decoration={markings, mark=at position 0.55 with {\arrow{Stealth}}}
    },
]

% --- Boucle électron-trou fermée : ellipse complète ---
% Centre (3,1.5), demi-axes a=1.2 (horizontal), b=1.0 (vertical)
% Vertex photon (cercles) aux points gauche (180°) et droit (0°) de l'ellipse
% Vertex phonon (carré) au point bas (-90°, soit (3,0.5))

% Branche supérieure gauche (électron) : de 180° à 0°, flèche au milieu
\draw[line, fermion] (1.8,1.5) arc (180:0:1.2 and 1.0);

% Branche inférieure droite (trou) : de 0° à -90°, flèche au milieu
\draw[line, fermion] (4.2,1.5) arc (0:-90:1.2 and 1.0);

% Branche inférieure gauche (trou) : de -90° à -180°, flèche au milieu
\draw[line, fermion] (3,0.5) arc (-90:-180:1.2 and 1.0);

% --- Vertex électron-photon (cercle), absorption -- point gauche (180°) ---
\draw[fill=white] (1.8,1.5) circle (0.18);

% --- Photon incident (ligne pointillée) ---
\draw[line, dashed] (0.2,1.5) -- (1.62,1.5);
\node[above] at (0.9,1.5) {$\omega_L$};

% --- Vertex électron-phonon (carré), point bas (-90°), milieu du propagateur ---
\draw[fill=white] (2.82,0.32) rectangle (3.18,0.68);

% --- Phonon sortant vers le bas depuis le vertex carré ---
\draw[line, decorate, decoration={snake, amplitude=0.8mm, segment length=3mm}]
    (3,0.32) -- (4,-0.7);
\node[below] at (4,-0.75) {$\omega_\nu$};

% --- Vertex électron-photon (cercle), émission -- point droit (0°), symétrique du premier ---
\draw[fill=white] (4.2,1.5) circle (0.18);

% --- Photon diffusé (ligne pointillée) ---
\draw[line, dashed] (4.38,1.5) -- (5.8,1.5);
\node[above] at (5.1,1.5) {$\omega_S$};

\end{tikzpicture}
\end{figure}